\documentclass[11pt]{article}
\usepackage[margin=1in]{geometry}
\usepackage{enumitem}
\setlist{nosep, before=\vspace{0.5em}, after=\vspace{0.5em}}
\usepackage{xcolor}
\usepackage{titling}
\usepackage{titlesec}
\usepackage[colorlinks=true,linkcolor=black,urlcolor=black,citecolor=black]{hyperref}

\titleformat{\section}{\large\bfseries}{\thesection.}{0.5em}{}
\titlespacing{\section}{0pt}{1.2em}{0.6em}
\titleformat{\subsection}{\normalsize\bfseries}{\thesubsection}{0.5em}{}
\titlespacing{\subsection}{0pt}{0.8em}{0.4em}
\titleformat{\subsubsection}{\normalsize\itshape}{\thesubsubsection}{0.5em}{}
\titlespacing{\subsubsection}{0pt}{0.6em}{0.3em}

\setlength{\droptitle}{-4em}
\pretitle{\begin{center}\large\bfseries}
\posttitle{\end{center}\vspace{-1em}}
\preauthor{\begin{center}\normalsize}
\postauthor{\end{center}\vspace{-1em}}
\predate{\begin{center}\normalsize}
\postdate{\end{center}\vspace{-0.5em}}

\title{Thesis Outline}
\author{Diego Toribio\\[0.3em]{\normalsize March 11, 2026}}
\date{\vspace{-2.5em}}

\begin{document}
\maketitle
\thispagestyle{empty}

\section{Introduction}

Establishes the data efficiency problem, identifies why existing
evaluations of self-supervised representation learning are
inconclusive, and states the research question.

\subsection{The Data Efficiency Problem}

Deep RL agents typically require hundreds of millions of training
interactions. The Atari-100K benchmark limits agents to 100K
interactions -- roughly two hours of real-time play -- exposing how
dependent these methods are on massive data.

\begin{itemize}[nosep]
\item The gap between human and agent data requirements
\item Why data efficiency matters beyond games: sim-to-real,
  costly environments, safety-critical domains
\item Four families of approaches: data augmentation, representation
  learning, world models, algorithmic improvements
\item This thesis focuses on the first two
\end{itemize}

\subsection{The Attribution Problem}

SPR adds a self-supervised objective that trains the encoder to predict
its own future representations. It achieved strong Atari-100K results,
but was only evaluated on top of Rainbow -- a complex agent with six
extensions that each change learning dynamics. We cannot attribute
SPR's gains to representation learning alone.

\begin{itemize}[nosep]
\item Rainbow's components (distributional values, prioritized replay,
  n-step returns, etc.) confound the result
\item Our early DQN+SPR results confirm this: only 1 of 6 games shows
  a strong lift, 4 games show flat or negative effects
\item This is not a failure of SPR -- it shows the base agent matters,
  and nobody has studied when or why
\end{itemize}

\subsection{Research Question and Approach}

\vspace{0.5em}
\begin{quote}
When does self-supervised representation learning improve
data-efficient reinforcement learning, and what determines whether
the learned representations are useful?
\end{quote}
\vspace{0.5em}

Three sub-questions:

\begin{itemize}[nosep]
\item \textbf{Isolation} -- Does SPR improve vanilla DQN? How does it
  interact with data augmentation? (each method alone and in
  combination, 6 games, 3 seeds)
\item \textbf{Capability} -- Does SPR need a capable base agent?
  (Rainbow+SPR on the same games as an experimental control)
\item \textbf{Interpretability} -- What do the representations encode,
  and why do they differ across capability levels? (linear probing,
  latent visualization, transition model analysis)
\end{itemize}

\subsection{Contributions}

\begin{itemize}[nosep]
\item \textbf{Isolation study} -- first evaluation of SPR on vanilla
  DQN, separating representation learning from algorithmic
  improvements
\item \textbf{Interaction study} -- testing augmentation and SPR alone
  and in combination to reveal whether they are independent or
  redundant
\item \textbf{Interpretability analysis} -- probing and visualization
  across capability levels, producing a finding about when and why
  self-supervised representations become useful
\item \textbf{Component attribution} (if investigation paths yield
  clean results) -- identifying which Rainbow component enables useful
  representations
\end{itemize}

\subsection{Thesis Outline}

\begin{itemize}[nosep]
\item Chapter 2 -- Background: RL, DQN, Rainbow, Atari-100K, SPR,
  augmentation, related work
\item Chapter 3 -- Methods: implementations, experiment design,
  interpretability methods
\item Chapter 4 -- Results: DQN isolation study, capability
  comparison, interpretability findings, cross-game generalization
\item Chapter 5 -- Discussion: interpretation, limitations, broader
  positioning
\item Chapter 6 -- Conclusion: contributions, implications, future work
\end{itemize}

\section{Background and Related Work}

Gives the reader the foundations needed to follow the experiment
design and interpret the results. Builds from RL basics through the
two agents that define the capability axis, then the benchmark, the
two methods under study, and positioning against the field.

\subsection{Reinforcement Learning Foundations}

\begin{itemize}[nosep]
\item Markov decision processes: states, actions, rewards, transitions,
  discount factor
\item Value functions and the Bellman equation
\item Q-learning and temporal difference updates
\item Why function approximation is needed for large state spaces
\end{itemize}

\subsection{Deep Q-Networks}

The ``simple agent'' in our comparison. This section covers what DQN
adds to tabular Q-learning and why it became the foundation for
deep RL.

\begin{itemize}[nosep]
\item Nature CNN architecture: convolutional layers over stacked frames
\item Experience replay: breaking correlation between consecutive samples
\item Target network: stabilizing learning with a frozen copy
\item Frame stacking and preprocessing pipeline
\end{itemize}

\subsection{Rainbow DQN}

The ``capable agent'' in our comparison -- not a contribution, but the
instrument that lets us test whether SPR needs a strong base agent.

\begin{itemize}[nosep]
\item Six extensions over vanilla DQN: distributional values (C51),
  prioritized replay, dueling architecture, n-step returns, noisy
  exploration, Double DQN
\item What each extension changes about learning dynamics
\item Why Rainbow became the standard base agent for SPR and other
  representation learning methods
\end{itemize}

\subsection{The Atari-100K Benchmark}

The testbed where both agents are evaluated. This section explains the
protocol and why it is the standard for data-efficient RL.

\begin{itemize}[nosep]
\item 100K environment steps vs the standard 200M -- what the constraint
  means in practice
\item Evaluation protocol: sticky actions, frame skip, evaluation
  schedule
\item Metrics: human-normalized score, IQM
\item Why Atari-100K rather than other benchmarks
\end{itemize}

\subsection{Self-Predictive Representations}

The representation learning method under study. This section explains
how SPR works and what questions remain about it.

\begin{itemize}[nosep]
\item Core idea: train the encoder to predict its own future
  representations
\item Architecture: transition model, EMA target encoder, projection
  and prediction heads
\item Training objective: cosine similarity loss over multi-step
  predictions
\item Originally evaluated only on Rainbow -- the attribution problem
  introduced in Chapter 1
\end{itemize}

\subsection{Data Augmentation for RL}

The second method under study. This section covers how augmentation
works as implicit regularization under limited data.

\begin{itemize}[nosep]
\item DrQ random shift: pad and crop game frames as a regularizer
\item Why augmentation helps in low-data regimes
\item The CURL finding: gains attributed to contrastive learning
  actually came from the augmentation
\end{itemize}

\subsection{Related Work and Positioning}

Where this thesis sits in the landscape of data-efficient RL methods.

\begin{itemize}[nosep]
\item World models: SimPLe, Dreamer -- learn a model of the environment
\item Planning methods: EfficientZero, MuZero -- combine model learning
  with search
\item Training efficiency: BBF, SR-SPR -- scale up existing methods
\item Why no prior work has isolated SPR from its base agent or compared
  representations across capability levels
\end{itemize}

\section{Methods}

Describes what we built, how the experiments are structured, and
how the interpretability analyses work. Enough detail that someone
could reproduce the setup.

\subsection{DQN Implementation}

All components implemented from scratch in PyTorch.

\begin{itemize}[nosep]
\item Nature CNN encoder: three convolutional layers over 4 stacked
  grayscale frames
\item Training procedure: RMSProp, epsilon-greedy exploration with
  linear decay, periodic target network updates
\item Replay buffer with circular storage
\item Atari preprocessing: frame skipping, grayscale conversion,
  84x84 resizing
\item Hyperparameter table and any deviations from Mnih 2015
\end{itemize}

\subsection{Data Augmentation}

DrQ-style random shift applied to sampled batches during training.

\begin{itemize}[nosep]
\item Pad frames to 92x92, random crop back to 84x84
\item Applied independently to both observations and next-observations
\item Configurable via YAML -- same toggle for DQN and Rainbow
  conditions
\end{itemize}

\subsection{Self-Predictive Representations}

SPR integrated identically on both DQN and Rainbow so the comparison
is controlled.

\begin{itemize}[nosep]
\item Transition model: two convolutional layers with action
  conditioning via one-hot broadcast
\item EMA target encoder with momentum 0.99
\item Projection and prediction heads (MLP)
\item Cosine similarity loss over K-step predicted representations
\item Combined loss: TD loss + weighted SPR loss
\item Encoder dropout (0.5) for regularization
\end{itemize}

\subsection{Rainbow Implementation}

Rainbow serves as infrastructure for the capability comparison.

\begin{itemize}[nosep]
\item Distributional value estimates (C51, 51 atoms)
\item Prioritized experience replay with importance sampling weights
\item Dueling network architecture
\item N-step returns (n=3)
\item Noisy exploration (factorised Gaussian)
\item Double DQN target computation
\item Per-game configs for all six games
\end{itemize}

\subsection{Experiment Design}

Two-part design testing SPR across two agents of different capability.

\begin{itemize}[nosep]
\item Part 1: DQN isolation study -- four conditions (baseline,
  +augmentation, +SPR, +both) across 6 games, 3 seeds each
\item Part 2: Rainbow comparison -- Rainbow and Rainbow+SPR on the
  same 6 games, 3 seeds each
\item Game selection rationale: six games spanning dense to moderate
  reward, chosen for diverse interaction patterns in published results
\item Evaluation: 30 episodes at final checkpoint, mean return with
  standard deviation
\item Compute: Google Colab Pro with automated Drive sync
\end{itemize}

\subsection{Interpretability Methods}

Three analyses on saved checkpoints, plus two investigation paths
for identifying which agent properties enable useful representations.

\begin{itemize}[nosep]
\item Linear probing: train classifiers on frozen encoder outputs to
  decode game-relevant features across all conditions
\item Latent space visualization: t-SNE/PCA of encoder outputs,
  comparing structure across DQN and Rainbow conditions
\item Transition model analysis: cosine similarity between predicted
  and actual next-states over K steps
\item Path B (correlational): correlate Rainbow training metrics with
  probe accuracy to narrow down which components matter
\item Path A (interventional): targeted ablations on the strongest
  candidates from Path B
\end{itemize}

% ================================================================
% CHAPTER 4: RESULTS
% ================================================================
% Question this chapter answers:
%   What happened when we ran the experiments?
%
% Minimal interpretation here -- present the data clearly.
% Discussion chapter handles the ``why.''
% ================================================================

\section{Results}

\subsection{DQN Isolation Study}
% Four conditions on vanilla DQN: baseline, +augmentation, +SPR,
% +both. Results per game, learning curves, and whether augmentation
% and SPR provide independent or redundant improvements.

\subsection{Rainbow Capability Comparison}
% Rainbow+SPR vs DQN+SPR across the same 6 games.
% The capability gap and its magnitude per game.

\subsection{Interpretability Analysis}
% Linear probe accuracy across conditions (DQN, DQN+SPR,
% Rainbow, Rainbow+SPR). Latent space visualizations.
% Transition model prediction quality comparison.
% Path B metric correlations. Path A ablation results.

\subsection{Cross-Game Generalization}
% Does the capability-representation relationship hold across
% dense, moderate, and sparse reward games?
% Game-specific vs universal findings.

% ================================================================
% CHAPTER 5: DISCUSSION
% ================================================================
% Question this chapter answers:
%   What do the results mean, why did they turn out this way,
%   and what are the limits of our claims?
% ================================================================

\section{Discussion}

\subsection{The Capability Threshold}
% What the DQN vs Rainbow comparison reveals about when
% self-supervised representation learning works. What the
% base agent provides that makes representations exploitable.

\subsection{Augmentation and Representation Learning}
% Are they independent or redundant? Why? Connection to
% CURL finding. Practical implications for practitioners.

\subsection{What Representations Encode}
% Synthesis of interpretability findings. What SPR learns
% on capable vs simple agents. Broader implications for
% understanding learned representations in RL.

\subsection{Limitations}
% Scale (6 games, 3 seeds), algorithm scope (DQN and Rainbow
% only), compute constraints, interpretability scope,
% augmentation scope (single type). Honest assessment of
% what the findings do and do not establish.

\subsection{Broader Context}
% Connection to self-supervised pre-training in NLP/vision.
% Positioning against published baselines (honest about not
% competing with SOTA). Implications for the field.

% ================================================================
% CHAPTER 6: CONCLUSION
% ================================================================
% Question this chapter answers:
%   What did we contribute, what does it imply, and what
%   should come next?
% ================================================================

\section{Conclusion}

\subsection{Summary of Contributions}
% Restate research question. Summarize each contribution
% with the actual findings (filled after experiments).

\subsection{Implications}
% For practitioners: when to use SPR, whether augmentation
% is worth adding, what to expect on simple agents.
% For researchers: what interpretability reveals, SPR
% isolation as a reference point for future work.

\subsection{Future Work}
% Concrete directions motivated by findings: expanded game
% set, additional DQN extensions, contrastive comparison,
% world model bridge, cross-game transfer.

\end{document}
